\documentclass[12pt,a4paper]{article}
\raggedbottom

\usepackage{ugmath}
\usepackage{float}
\usepackage{placeins}
\usepackage[labelfont=bf,labelsep=space]{caption}
\usepackage{pdfpages}
\usepackage{tikz-cd}
\usepackage{amsthm}
\renewcommand{\Tr}{\operatorname{Tr}}
\newcommand{\llaves}[1]{\{#1\}}
\newcommand{\parentesis}[1]{\left(#1\right)}
\newcommand{\corchetes}[1]{\left[#1\right]}
\newcommand{\llavesgrande}[1]{\big\{#1\big\}}
\newcommand{\llavesgrandes}[1]{\left\{#1\right\}}


\begin{document}
        
    \thispagestyle{empty}
    \includepdf[pages=1]{main.pdf}

    \begin{problema}
        Prueba los incisos (b) y (c) de la proposición 3.
    \end{problema}

    \begin{proof}
        \begin{enumerate}
            \item[(b)] Si $x\in A\cup B$ entonces $x\in A$ y como $A\subseteq B$ se tiene que $x\in B$ y, por lo tanto
            $A\cup B\subseteq B$. Ahora si $x\in B$ entonces como $A\subseteq B$, obtenemos $x\in A$ y por definición de unión
            $x\in A\cup B$, por consiguiente $B\subseteq A\cup B$. Concluimos que,
            \begin{equation*}
                A\cup B=B.
            \end{equation*}
            Si $x\in A\cap B$ entonces por definición de intersección $x\in A$ y, por tanto $A\cap B\subseteq A$. Ahora si $x\in A$ entonces como $A\subseteq B$
            tenemos que $x\in B$ y por definición de intersección $x\in A\cap B$, y asi, $A\subseteq A\cap B$.
            Concluimos que
            \begin{equation*}
                A\cap B=A.
            \end{equation*}
            Si $x\in A\setminus B$ entonces $x\in A\wedge x\notin B$, pero esto es imposible porque $A\subseteq B$, por lo que no existe $x$
            tal que $x\in A\setminus B$. En consecuencia, $A\setminus B$ no tiene elementos por ende,
            \begin{equation*}
                A\setminus B=\varnothing.
            \end{equation*}
            \item[(c)] Si $x\in A\cup C$ entonces $x\in A$, pero como $A\subseteq B$ entonces $x\in B$ y por lo tanto
            $x\in B\cup C$. Concuimos que
            \begin{equation*}
                A\cup C\subseteq B\cup C
            \end{equation*}
            Si $x\in A\cap C$ entonces $x\in A$ y como $A\subseteq B$ entonces $x\in B$ y por lo tanto $x\in B\cap C$. Concluimos que
            \begin{equation*}
                A\cap C\subseteq B\cap C
            \end{equation*}
            Si $x\in C\setminus B$ entonces $x\in C\wedge x\notin B$, como $A\subseteq B$ entonces si $x\notin B$ se tiene $x\notin A$
            De esta forma como $x\in C\setminus B$ tambien se tiene que $x\in C$ y por lo tanto $x\in C\wedge x\notin A$. Esto es,
            \begin{equation*}
                C\setminus B\subseteq C\setminus A.
            \end{equation*}
        \end{enumerate}
    \end{proof}

    \begin{problema}
        Prueba el inciso (b) de la proposición 5.
    \end{problema}

    \begin{proof}
        Supongamos que existe $y\in\R$ y $x\neq0$ tal que
        \begin{equation*}
            x\cdot y=x.
        \end{equation*}
        Por la regla de cancelación de producto obtenemos,
        \begin{equation*}
            y=x\cdot x^{-1}.
        \end{equation*}
        Así, por la unicidad del inverso multiplicativo
        \begin{equation*}
            y=1.
        \end{equation*}
        Por lo que el único número $y$ tal que cumple con la expresión inicial es 1. Por lo tanto existe un unico nuetro multiplicativo y es el 1.
    \end{proof}
    
    \begin{problema}
        \textbf{(Leyes de los signos para la multiplicaci n)} Prueba que
        si $x$ y $y$ son n meros reales, entonces
        \begin{enumerate}
            \item[(a)] $(+x)\left(+y\right)=+(xy).$ 
            \item[(b)] $(-x)\left(+y\right)=-(xy).$
            \item[(c)] $\left(+x\right)\left(-y\right)=-\left(xy\right).$
            \item[(d)] $(-x)(-y)=+\left(xy\right).$
        \end{enumerate}
        Las leyes de los signos para la multiplicaci n se resumen de manera
        conveniente como 
        \begin{equation*}
            \begin{array}{ccccc}
            + & \text{por} & + & \text{da} & +\\
            - & \text{por} & + & \text{da} & -\\
            + & \text{por} & - & \text{da} & -\\
            - & \text{por} & - & \text{da} & +
            \end{array}
        \end{equation*}
        \textbf{Sugerencia:} El inciso (a) es claro, pues $+x=x,$ $+y=y$
        y $+\left(xy\right)=xy.$ Usa la unicidad del inverso aditivo para
        probar los incisos (b) y (c). Finalmente, deduce (d) de los incisos
        (b) y (c).
    \end{problema}

    \begin{proof}
        \begin{enumerate}
            \item[(a)] $(+x)(+y)=+(xy)$.\\
            Es claro pues $+x=x$, $+y=y$ y $+(xy)$.
            \item[(b)] $(-x)(+y)=-(xy)$.\\
            Consideremos $(-x)y$. Escribimos $-x$ como $(-1)\cdot x$. Entonces,
            \begin{equation*}
                (-x)y=[(-1)x]y.
            \end{equation*}
            Por la asociatividad de el producto,
            \begin{equation*}
                [(-1)\cdot x]\cdot y=(-1)(x\cdot y).
            \end{equation*}
            Ahora por el problema 4 de la tarea 2, se cumple que
            \begin{equation*}
                (-1)(xy)=-(xy).
            \end{equation*}
            Por lo tanto, $(-x)(+y)=-(xy)$.
            \item[(c)] $(+x)(-y)=-(xy)$.\\
            Consideremos $(-y)(x)$. Escribimos $-y$ como $(-1)\cdot y$. Entonces,
            \begin{equation*}
                (x)(-y)=x\cdot[(-1)\cdot y].
            \end{equation*}
            Por la asociatividad de la multiplicación,
            \begin{equation*}
                x\cdot[(-1)\cdot y]=(-1)(xy).
            \end{equation*}
            Así, nuevamente por el problema 4 de la tarea 2 se cumple que
            \begin{equation*}
                (-1)(xy)=-(xy).
            \end{equation*}
            \item[(d)] $(-x)(-y)=(+xy)$.\\
            Escribimos los números negativos en términos de $-1$,
            \begin{equation*}
                -x=(-1)\cdot x,\qquad -y=(-1)\cdot y.
            \end{equation*}
            Entonces su producto es
            \begin{equation*}
                (-x)(-y)=[(-1)\cdot x]\cdot[(-1)\cdot y].
            \end{equation*}
            Aplicando la asociatividad y la conmutatividad tenemos que
            \begin{equation}
                (-x)(-y)=[(-1)\cdot x]\cdot[(-1)\cdot y]=[(-1)\cdot(-1)]\cdot(xy)
            \end{equation}
            Ahora calculamos $(-1)\cdot(-1)$. Sabemos que $0=0\cdot0$, pero $0=(1+(-1))$. Entonces,
            \begin{equation*}
                0+0=(1+(-1))\cdot(1+(-1)).
            \end{equation*}
            Expandiendo usando la distributividad,
            \begin{align*}
                (1+(-1))\cdot(1+(-1))&=(1+(-1))\cdot1+(1+(-1))\cdot(-1)\\
                &=1\cdot1+(-1)\cdot1+1\cdot(-1)+(-1)\cdot(-1)\\
                &=1+(-1)+(-1)+(-1)\cdot(-1)\\
                &=(1-1-1)+(-1)(-1)\\
                &=-1+(-1)(-1).
            \end{align*}
            Como $0\cdot0=0$ tenemos que
            \begin{equation*}
                0=-1+(-1)(-1)
            \end{equation*}
            de aquí,
            \begin{equation*}
                (-1)\cdot(-1)=1
            \end{equation*}
            sustituyendo en (1)
            \begin{equation*}
                (-x)(-y)=[(-1)(-1)](xy)=1\cdot(xy)=xy.
            \end{equation*}
            Así,
            \begin{equation*}
                (-x)(-y)=-(xy).
            \end{equation*}
        \end{enumerate}
    \end{proof}

    \begin{problema}
        Sean $x$ y $y$ números reales. Prueba que si $xy=0$, entonces $x=0$ o $y=0$. \textbf{Segerencia:} Procede por contradicción.
    \end{problema}

    \begin{proof}
        Supongamos que $xy=0$. Procedamos por contradicción, supongamos que
        \begin{equation*}
            x\neq0\qquad\text{y}\qquad y\neq0.
        \end{equation*}
        Como $x\neq0$ entonces existe $x^{-1}\in\R$ tal que
        \begin{equation*}
            x^{-1}\cdot x=1.
        \end{equation*}
        Multiplicando la igualdad $xy=0$ por $x^{-1}$, obtenemos
        \begin{equation*}
            x^{-1}xy=x^{-1}0,
        \end{equation*}
        y, por la asociatividad, la unicidad del inverso multiplicativo, se tiene que
        \begin{equation*}
            y=0
        \end{equation*}
        pero esto contradice que $y\neq0$. Por tanto, concluimos que
        \begin{equation*}
            x=0\text{ ó } y=0
        \end{equation*}
    \end{proof}
    
    \begin{problema}
        Prueba que si $0\leq x\lt y$ y $0\leq z\lt w$, entonces $xz\lt yw$.
    \end{problema}

    \begin{proof}
        Por hipótesis tenemos que
        \begin{equation*}
            x\lt y\qquad\text{ y }\qquad z\lt w.
        \end{equation*}
        De la primer desigualdad, dado que $0\leq x\lt y$, entonces, por la preservación del orden bajo el producto por numeros positivosm, se tiene
        \begin{equation*}
            xz\lt yz.
        \end{equation*}
        Por otra parte, de la segunda desigualdad, nuevamente por la preservación del orden bajo el producto por numeros positivos, se tiene
        \begin{equation*}
            yz\lt yw.
        \end{equation*}
        Así, por transitividad se concluye que
        \begin{equation*}
            xz\lt yw.
        \end{equation*}
    \end{proof}

    \begin{problema}
        Prueba la proposición 12.
    \end{problema}

    \begin{proof}
        Como $A$ es acotado inferiormente, existe $m\in\mathbb{R}$ tal que
        \begin{equation*}
            m\leq a,\qquad\forall a\in A.
        \end{equation*}
            Consideremos el conjunto
        \begin{equation*}
            -A={-a:a\in A}.
        \end{equation*}
        Primero, $-A\neq\varnothing$, pues $A\neq\varnothing$. Además, $-A$ está acotado superiormente. En efecto, si $m$ es una cota inferior de $A$, entonces
        \begin{equation*}
            m\leq a,\qquad\forall a\in A,
        \end{equation*}
        y multiplicando por $-1$, invirtiendo la desigualdad,
        \begin{equation*}
            -a\leq-m,\qquad\forall a\in A.
        \end{equation*}
        Por tanto, $-m$ es una cota superior de $-A$. Por el axioma del supremo, existe
        \begin{equation*}
            s=\sup(-A).
        \end{equation*}
        Definimos
        \begin{equation*}
            \alpha=-s.
        \end{equation*}
        Afirmamos que
        \begin{equation*}
            \alpha=\inf(A).
        \end{equation*}
        Veamos que cumple las dos propiedades de la definición de ínfimo. Primero, como
        \begin{equation*}
            s=\sup(-A),
        \end{equation*}
        tenemos
        \begin{equation*}
            -a\leq s,\qquad\forall a\in A.
        \end{equation*}
        Multiplicando por $-1$, invirtiendo la desigualdad,
        \begin{equation*}
            -s\leq a,\qquad\forall a\in A.
        \end{equation*}
        Es decir,
        \begin{equation*}
            \alpha\leq a,\qquad\forall a\in A.
        \end{equation*}
        Luego $\alpha$ es una cota inferior de $A$. Sea ahora $l$ una cota inferior cualquiera de $A$. Entonces
        \begin{equation*}
            l\leq a,\qquad\forall a\in A.
        \end{equation*}
        Por tanto,
        \begin{equation*}
            -a\leq-l,\qquad\forall a\in A,
        \end{equation*}
        de modo que $-l$ es una cota superior de $-A$. Como
        \begin{equation*}
            s=\sup(-A),
        \end{equation*}
        necesariamente
        \begin{equation*}
            s\leq-l.
        \end{equation*}
        Multiplicando por $-1$, invirtiendo la desigualdad,
        \begin{equation*}
            -s\geq l.
        \end{equation*}
        Por la definición de $\alpha$,
        \begin{equation*}
            l\leq\alpha.
        \end{equation*}
        Por tanto, $\alpha$ es una cota inferior de $A$ mayor o igual que cualquier otra cota inferior de $A$. Luego,
        \begin{equation*}
            \alpha=\inf(A).
        \end{equation*}
    \end{proof}

    \begin{problema}
        Prueba que para todo número natural $n$ se cumple que
        \begin{equation*}
            n\lt n+1.
        \end{equation*}
    \end{problema}

    \begin{proof}
        Sea $n=1$. Entonces
        \begin{equation*}
            1\lt 2.
        \end{equation*}
        Supongamos ahora que para algun $n\in\N$, se cumple
        \begin{equation*}
            n\lt n+1.
        \end{equation*}
        Debemos demostrar que también es verdadera para $n+1$, es decir, que
        \begin{equation*}
            n+1\lt n+2
        \end{equation*}
        Por la hipotesis de inducción sabemos que
        \begin{equation*}
            n\lt n+1.
        \end{equation*}
        Por tanto,
        \begin{equation*}
            n+1\lt (n+1)+1.
        \end{equation*}
        Esto es,
        \begin{equation*}
            n+1\lt n+2.
        \end{equation*}
        Y esto es precisamente lo que queríamos demostrar.
    \end{proof}

    \begin{problema}
        Determina si las reglas de correspondencia $f$ de los cuadros de
        los incisos siguientes son funciones (argumenta tu respuesta).
        \begin{center}
            (a)~~~%
            \begin{tabular}{|c|c|}
                \hline 
                \hspace{0.2cm}$x$\hspace{0.2cm} & $f\left(x\right)$\tabularnewline
                \hline 
                \hline 
                $1$ & $0$\tabularnewline
                \hline 
                $2$ & $5$\tabularnewline
                \hline 
                $3$ & $7$\tabularnewline
                \hline 
                $4$ & $7$\tabularnewline
                \hline 
                $5$ & $8$\tabularnewline
                \hline 
            \end{tabular}~~~~~~(b)~~~%
            \begin{tabular}{|c|c|}
                \hline 
                \hspace{0.2cm}$x$\hspace{0.2cm} & $f\left(x\right)$\tabularnewline
                \hline 
                \hline 
                $1$ & 0\tabularnewline
                \hline 
                $2$ & 0\tabularnewline
                \hline 
                $3$ & 0\tabularnewline
                \hline 
                $4$ & 0\tabularnewline
                \hline 
                $5$ & 0\tabularnewline
                \hline 
            \end{tabular}~~~~~~(c)~~~%
            \begin{tabular}{|c|c|}
                \hline 
                \hspace{0.2cm}$x$\hspace{0.2cm} & $f\left(x\right)$\tabularnewline
                \hline 
                \hline 
                $1$ & $3$\tabularnewline
                \hline 
                $2$ & $3$\tabularnewline
                \hline 
                $3$ & $5$\tabularnewline
                \hline 
                $4$ & $7$\tabularnewline
                \hline 
                $4$ & $6$\tabularnewline
                \hline 
            \end{tabular}\par
        \end{center}
    \end{problema}

    \begin{solucion}
        \begin{enumerate}
            \item[(a)] Es una función ya que le asigna a cada elemento $x$ un único elemento de $f(x)$.
            \item[(b)] Es una función ya que le asigna a cada elemento $x$ un único elemento de $f(x)$.
            \item[(c)] No es una función ya que le asigna más de un elemento de $x$ a un único elemento de $f(x)$.  
        \end{enumerate}
    \end{solucion}

    \begin{problema}
        Calcula el dominio de las funciones siguientes:
        \begin{multicols}{2}
            \begin{enumerate}
                \item[(a)] $f(x)=x^{2}-5.$
                \item[(b)] $f(x)=\sqrt{x+2}.$
                \item[(c)] $f(x)=\frac{2}{x-3}.$
            \end{enumerate}
        \end{multicols}
    \end{problema}

    \begin{solucion}
        \begin{enumerate}
            \item[(a)] $f(x)=x^{2}-5$.\\
            Ya que es un polinomio y por lo tanto no tiene mas restricciones el dominio de $f$ es $\R$.
            \item[(b)] $f(x)=\sqrt{x+2}$.\\
            Dado que la raíz cuadrada de números negativos no esta definida en $\R$, el domino de la función $f$ es
            \begin{equation*}
                \{x:x\geq-2\}=[-2,\infty).
            \end{equation*}
            \item[(c)] $f(x)=\frac{2}{x-3}$.\\
            Dado que la división por 0 no esta definida, el dominio de la función $f$ es
            \begin{equation*}
                \R\setminus\{3\}=(-\infty,3)\cup(3,+\infty).
            \end{equation*}
        \end{enumerate}
    \end{solucion}

\end{document}